\documentclass[11pt]{article}

\usepackage{amsmath, amssymb, bm}
\usepackage{geometry}
\usepackage{hyperref}
\usepackage{graphicx}
\geometry{margin=2.5cm}

\title{Technical Design and Software Architecture of the RHF Core in \texttt{QuetzalQC}}
\author{}
\date{}

\begin{document}

\maketitle

\section{Overview}

\texttt{QuetzalQC} is designed as a modular quantum chemistry research engine.
The current implementation focuses on a Restricted Hartree--Fock (RHF) solver
built using PyTorch for linear algebra operations, with integral evaluation
delegated to a backend provider (PySCF).

The architecture separates:

\begin{itemize}
\item Mathematical formulation
\item Linear algebra utilities
\item Density construction
\item Coulomb and exchange contraction
\item SCF control flow
\item Backend integral provider
\end{itemize}

This modularization allows independent optimization and future extensibility.

\section{Module Architecture}

The core structure is:

\begin{verbatim}
quetzalqc/
  backends/
      pyscf_backend.py
  core/
      linalg.py
      density.py
      jk.py
      energy.py
      diis.py
      rhf.py
\end{verbatim}

Each module corresponds directly to a mathematical component of the RHF equations.

\section{Backend Abstraction Layer}

\subsection{PySCF Backend}

The backend provides:

\begin{itemize}
\item Overlap matrix $\mathbf{S}$
\item One-electron Hamiltonian $\mathbf{H}$
\item Two-electron integrals $(\mu\nu|\lambda\sigma)$
\item Nuclear repulsion energy $E_{nuc}$
\item Electron count $N$
\end{itemize}

The SCF driver is backend-agnostic; only these tensors are required.
This enables replacement by:

\begin{itemize}
\item Custom integral engines
\item Density-fitting backends
\item GPU-based integral evaluation
\end{itemize}

\section{Linear Algebra Layer}

File: \texttt{core/linalg.py}

Responsible for solving the generalized eigenvalue problem:

\[
\mathbf{F}\mathbf{C} = \mathbf{S}\mathbf{C}\bm{\varepsilon}.
\]

The implementation uses symmetric orthogonalization:

\[
\mathbf{X} = \mathbf{S}^{-1/2}.
\]

Transformation:

\[
\mathbf{F}' = \mathbf{X}^T \mathbf{F} \mathbf{X}.
\]

This reduces the problem to a standard Hermitian diagonalization.

This design ensures:

\begin{itemize}
\item Numerical stability
\item Compatibility with PyTorch eigensolvers
\item Clear separation of mathematical transformation
\end{itemize}

\section{Density Module}

File: \texttt{core/density.py}

Implements:

\[
\mathbf{P} = 2 \mathbf{C}_{occ} \mathbf{C}_{occ}^T.
\]

This is independent of how orbitals were obtained,
allowing future reuse for:

\begin{itemize}
\item DFT
\item Post-HF methods
\item Finite-temperature occupations
\end{itemize}

\section{Coulomb and Exchange Contractions}

File: \texttt{core/jk.py}

Implements tensor contractions:

\[
J_{\mu\nu} =
\sum_{\lambda\sigma}
P_{\lambda\sigma}
(\mu\nu|\lambda\sigma),
\]

\[
K_{\mu\nu} =
\sum_{\lambda\sigma}
P_{\lambda\sigma}
(\mu\lambda|\nu\sigma).
\]

Implemented using \texttt{torch.einsum}.

This module isolates the $O(N^4)$ scaling component,
facilitating later optimization via:

\begin{itemize}
\item Density fitting
\item GPU kernels
\item CUTLASS tensor contractions
\end{itemize}

\section{Energy Evaluation}

File: \texttt{core/energy.py}

Implements the RHF total energy:

\[
E =
\frac{1}{2}
\mathrm{Tr}
\left[
\mathbf{P}
(\mathbf{H} + \mathbf{F})
\right]
+ E_{nuc}.
\]

Separated to allow:

\begin{itemize}
\item Alternative energy expressions
\item DFT energy functional insertion
\item Energy decomposition analysis
\end{itemize}

\section{DIIS Acceleration}

File: \texttt{core/diis.py}

Implements Direct Inversion in the Iterative Subspace (DIIS),
using the commutator error:

\[
\mathbf{e} = \mathbf{F}\mathbf{P}\mathbf{S}
- \mathbf{S}\mathbf{P}\mathbf{F}.
\]

This improves convergence robustness without modifying
core RHF equations.

\section{SCF Driver}

File: \texttt{core/rhf.py}

The SCF driver orchestrates:

\begin{enumerate}
\item Initial guess from core Hamiltonian
\item Fock construction
\item Generalized eigenvalue solve
\item Density update
\item Convergence test
\end{enumerate}

This driver contains no integral logic,
no hardcoded basis assumptions,
and no backend dependency beyond tensor inputs.

\section{Data Flow}

The RHF iteration cycle is:

\[
\mathbf{P}^{(n)}
\rightarrow
\mathbf{F}^{(n)}
\rightarrow
\mathbf{C}^{(n)}
\rightarrow
\mathbf{P}^{(n+1)}.
\]

This defines a nonlinear fixed-point problem:

\[
\mathbf{P} = \mathcal{F}(\mathbf{P}).
\]

\section{Design Principles}

\begin{itemize}
\item \textbf{Modularity:} Each mathematical concept corresponds to a module.
\item \textbf{Backend abstraction:} Integrals separated from SCF logic.
\item \textbf{Numerical transparency:} Explicit tensor operations.
\item \textbf{Extensibility:} Ready for DFT, post-HF, and GPU acceleration.
\item \textbf{Reproducibility:} Validated against PySCF.
\end{itemize}

\section{Future Extensions}

Planned extensions include:

\begin{itemize}
\item LDA/GGA DFT modules
\item Density fitting (RI) acceleration
\item CUDA-optimized tensor contractions
\item Analytical nuclear gradients
\item Born--Oppenheimer molecular dynamics
\end{itemize}

\end{document}

